\begin{helper}[Two disjoint exit paths]
Let $G=(V,E)$ be a finite loopless multigraph, let $T$ be a
tree-sequence, and let $\varphi$ be a proper partial $k$-edge-coloring.
Suppose that $V(T)$ is elementary under $\varphi$.  Let $a,b\in[k]$ be
distinct, suppose that the singleton $\{a\}$ is closed on $T$, and
suppose that $a,b$ are not $T$-interchangeable.  Then there are vertices
$a_1,b_1,a_2,b_2$ and two walks $P_i$ in the $(a,b)$-incidence graph
from $a_i$ to $b_i$ such that both $P_i$ are $T$-exit paths,
\[
  \operatorname{supp}(P_1)\cap\operatorname{supp}(P_2)=\varnothing,
  \qquad a_1\ne a_2,\qquad b_1\ne b_2.
\]
No assumption that $a$ is missing on $T$ is made.
\end{helper}
\begin{proof}
Write $X=V(T)$.  By the negation of \noderef{TreeInterchangeable}, there
are nodes $z_1,z_2$ whose \noderef{TwoColorComponent}s $C_1,C_2$ are
distinct \noderef{TwoColorPathComponent}s and each contain a vertex-node
from $X$.  By \noderef{TwoColorComponentsDisjoint}, the two components
are disjoint, including at exact edge-nodes.

We first describe any one of these components.  If it contains an
edge-node, it is a \noderef{TwoColorNontrivialPathComponent}.  The
spanning simple path given by \noderef{TwoColorComponentSpanningPath}
uses every node of the component.  Every incidence-graph node has degree
at most two by properness and
\noderef{TwoColorIncidenceDegreeAtMostTwo}.  The path-component
certificate supplies a node with at most one neighbor.  Consequently the
component has no extra adjacency joining the ends and no extra adjacency
from an end to an internal node: either kind of adjacency would make the
first kind of node have degree two everywhere, or give one of the
vertices involved degree at least three.  Thus the component itself is a
nontrivial simple path, and its two ends have degree one.

Each end is a vertex-node.  Indeed, the neighbors of a colored edge-node
in \noderef{TwoColorIncidenceGraph} are its incident vertex-nodes.
Incidence is membership in its unordered endpoint pair by
\noderef{Incident}; the loopless law in \noderef{LooplessMultigraph}
makes those endpoints distinct.  Hence a colored edge-node has two
neighbors.  Each vertex end misses $a$ or $b$: otherwise
\noderef{PresentColors} supplies distinct incident edges of colors $a$
and $b$, giving two distinct neighbors, contrary to degree one.  The
translation from absence among present colors to missingness is exactly
\noderef{MissingColors}.

If a component meeting $X$ contains no edge-node, it is the singleton
$\{\operatorname{inl}(x)\}$ for some $x\in X$.  Any incident edge of
color $a$ or $b$ would be an adjacent edge-node in the same component,
so both colors are missing at $x$ by \noderef{PresentColors} and
\noderef{MissingColors}.  There cannot be two such singleton components:
their distinct tree vertices would both miss $a$, contrary to
\noderef{ColoredVertexSetElementary}.

We next choose two ends outside $X$.  If $b$ is not missing on $T$, no
component is a singleton.  Among the four pairwise distinct ends of the
two disjoint nontrivial paths, an end in $X$ cannot miss $b$ and hence
misses $a$.  Elementariness permits at most one such end, so at least
three ends are outside $X$.  If $b$ is missing on $T$ and both components
are nontrivial, elementariness permits at most one of their four ends in
$X$ to miss $a$ and at most one to miss $b$; at least two ends are again
outside.  Finally, if one component is the singleton at $x\in X$, the
other is nontrivial and neither of its ends can be in $X$, since $x$
already misses both colors.  Thus in every case there are two distinct
outside ends on nontrivial component paths.

Orient each selected path from its selected outside end $a_i$ and stop
at its first vertex-node $b_i$ belonging to $X$.  This first position
exists because the component meets $X$; take the least such position on
the finite path.  The resulting prefix $P_i$ is simple, begins outside
$X$, ends in $X$, and has no other vertex-node in $X$.  Its outside end
misses $a$ or $b$.  Since both endpoints are vertex-nodes in the
bipartite incidence graph and are distinct, the prefix has length at
least two.  These are precisely the clauses of \noderef{TreeExitPath}.

If the chosen ends lie in different components, disjointness of the
components gives disjoint supports and distinct tree ends.  Otherwise
the chosen ends are the opposite ends of one simple path.  In the order
from one selected end to the other, the first prefix ends at the first
encounter with $X$, while the reversed prefix ends at the last encounter
with $X$.  The prefixes can meet only if those encounters are the same
vertex $v$.  In that event, let $f_1,f_2$ be the exact edge-nodes directly
before $v$ on the two prefixes.  Simplicity makes $f_1\ne f_2$.  The
preceding vertex on each prefix is outside $X$, whereas $v\in X$.
The adjacency description in \noderef{TwoColorIncidenceGraph}, together
with \noderef{Incident}, \noderef{LooplessMultigraph}, and
\noderef{BoundaryEdges}, therefore makes each $f_i$ a boundary edge of
$X$ whose color is $a$ or $b$.  Closure of $\{a\}$, as defined by
\noderef{TreeColorSetClosed} and \noderef{BoundaryColorEdges}, excludes
color $a$, so both edges have color $b$.  They are distinct and incident
with $v$, hence adjacent by \noderef{AdjacentEdges}; this contradicts
properness in \noderef{IsProperEdgeColoring}.  The exceptional overlap is
impossible.  Therefore the two selected exit paths have disjoint support,
and their outside and tree ends are respectively distinct, as claimed.
\end{proof}
